% ============================================================
% preamble.tex — shared style for the Schemes self-study reader
% Engine: XeLaTeX / tectonic.  Palatino-style via TeX Gyre Pagella.
% ============================================================
\usepackage{fontspec}
\usepackage{unicode-math}
% Use the tectonic-bundled defaults (Latin Modern + Latin Modern Math) — the
% classic TeX book look, guaranteed available without system font lookup.

\usepackage[a4paper,marginparwidth=2.4cm,marginparsep=0.35cm,
            inner=2.6cm,outer=4.6cm,top=2.6cm,bottom=2.8cm]{geometry}
\usepackage{microtype}
\usepackage{amsmath,amsthm,mathtools}
\usepackage{xcolor}
\usepackage{tikz}
\usetikzlibrary{arrows.meta,positioning,calc,decorations.markings,backgrounds}
\usepackage{tikz-cd}
\usepackage[most]{tcolorbox}
\usepackage{booktabs}
\usepackage{array}
\usepackage{enumitem}
\usepackage{caption}
\captionsetup{font=small,labelfont={bf,color=accentB}}
\usepackage{marginnote}
\usepackage{hyperref}

% ---- palette (matches the HTML course) ----
\definecolor{ink}{HTML}{1A1814}
\definecolor{inksoft}{HTML}{4A4640}
\definecolor{accent}{HTML}{7C1D28}   % burgundy
\definecolor{accentB}{HTML}{1F4E5F}  % teal
\definecolor{algc}{HTML}{5B3A86}     % violet
\definecolor{good}{HTML}{2F6B3A}
\definecolor{bad}{HTML}{9A2B2B}
\definecolor{warnc}{HTML}{B5651D}
\definecolor{paper2}{HTML}{F5F1E8}
\definecolor{rulec}{HTML}{D8D2C4}

\hypersetup{colorlinks=true,linkcolor=accentB,urlcolor=accent,citecolor=accent,
            pdftitle={Schemes: a self-study reader},pdfauthor={Schemes self-study}}
\color{ink}

% ---- math shortcuts ----
\newcommand{\Spec}{\operatorname{Spec}}
\newcommand{\Frac}{\operatorname{Frac}}
\newcommand{\colim}{\operatorname*{colim}}
\newcommand{\Sh}{\operatorname{Sh}}
\newcommand{\PSh}{\operatorname{PSh}}
\renewcommand{\AA}{\mathbb{A}}
\newcommand{\ZZ}{\mathbb{Z}}
\newcommand{\QQ}{\mathbb{Q}}
\newcommand{\RR}{\mathbb{R}}
\newcommand{\CC}{\mathbb{C}}
\newcommand{\FF}{\mathbb{F}}
\renewcommand{\O}{\mathcal{O}}
\newcommand{\sF}{\mathcal{F}}
\newcommand{\fp}{\mathfrak{p}}
\newcommand{\fq}{\mathfrak{q}}
\newcommand{\fm}{\mathfrak{m}}
\newcommand{\res}{\kappa}                       % residue field kappa
\newcommand{\restr}[1]{\!\restriction_{#1}}

% ---- Stacks Project tag + citation helpers ----
\newcommand{\Tag}[1]{\href{https://stacks.math.columbia.edu/tag/#1}{\texttt{#1}}}
\newcommand{\scite}[2]{\footnote{#2~(Stacks \Tag{#1}).}}
\newcommand{\alg}[1]{\textcolor{algc}{#1}}
\newcommand{\geo}[1]{\textcolor{accentB}{#1}}
\newcommand{\term}[1]{\textcolor{accent}{\itshape #1}}

% ---- callout boxes ----
\newtcolorbox{defbox}[1][Definition]{%
  enhanced,breakable,colback=paper2,colframe=accent,
  boxrule=0pt,leftrule=3pt,arc=2pt,left=8pt,right=8pt,top=6pt,bottom=6pt,
  fonttitle=\bfseries\scshape\color{accent},title={#1}}
\newtcolorbox{intubox}[1][Intuition]{%
  enhanced,breakable,colback=accentB!6,colframe=accentB,
  boxrule=0pt,leftrule=3pt,arc=2pt,left=8pt,right=8pt,top=6pt,bottom=6pt,
  fonttitle=\bfseries\scshape\color{accentB},title={#1}}
\newtcolorbox{warnbox}[1][Warning]{%
  enhanced,breakable,colback=warnc!8,colframe=warnc,
  boxrule=0pt,leftrule=3pt,arc=2pt,left=8pt,right=8pt,top=6pt,bottom=6pt,
  fonttitle=\bfseries\scshape\color{warnc},title={#1}}
\newtcolorbox{primarybox}{%
  enhanced,breakable,colback=paper2,colframe=accentB,boxrule=0.6pt,
  sharp corners,arc=2pt,left=8pt,right=8pt,top=6pt,bottom=6pt,
  borderline={0.6pt}{0pt}{accentB,dashed}}

% ---- retrieval-practice quiz item ----
% \quizitem{stem}{A}{B}{C}{D}{answer letter}{why}
\newcommand{\quizitem}[7]{%
  \item #1\par\smallskip
  \textbf{(A)} #2\quad\textbf{(B)} #3\quad\textbf{(C)} #4\quad\textbf{(D)} #5\par\smallskip
  {\small\color{good!75!black}\textbf{Answer: #6.}~#7}\par\medskip}
\newenvironment{quizlist}[1][Retrieval check]%
  {\par\smallskip\noindent\textbf{\scshape\color{accentB}#1}\begin{enumerate}[leftmargin=1.4em,itemsep=2pt]}%
  {\end{enumerate}}

% ---- dictionary / worked-example tables use booktabs directly ----
\newcolumntype{L}[1]{>{\raggedright\arraybackslash}p{#1}}

\theoremstyle{plain}
\newtheorem{theoremx}{Theorem}[chapter]
